\section{Proof}

\subsection{Boxes form a semi-ring}

\begin{proposition}\label{prop:RHsemiring}
For fixed $x$, the right handed rectangles $B_{x,a,b,c,d}$ form a semi-ring.
\end{proposition}
\begin{proof}
Show finite additivity by using whiskers argument. Idea:
 Induction on the number of vertical+horizontal
coordinates that appear. Look at any vertical division that reaches
the upper boundary. If there isn't one, there's a rectangle that
reaches clear across the top; remove that rectangle, by induction
build up the rest using side-by-side additivity, and then add that
missing rectangle.

Otherwise, divide the entire collection of rectangles at the location
of that vertical division. (This is possible, since this is the
correct direction for the whisker.) By induction, build up each half
by side-by-side additivity, and then attach the two sides together.
\end{proof}

From now on, let's denote the semiring by $\mathcal{S}$.

\subsection{$\mu_{\nu}$ is a measure}

This is the heart of the proof. Here we show that all the necessary hypothesis are satisfied to apply Theorem~\ref{thm:caratheodory}.

\subsubsection{Preliminar lemmas}

To prove useful facts about this $\mu_{\nu}$, we need some preliminar lemmas.

\begin{lemma}\label{lem:join-split-geom}
  If $a$ and $b$ cross $x$ to the right, then there is an integer $k$
  so that, for all $n,m,r$ with either $0 < r < m-n+k$ or $0 > r > m-n+k$,
  \begin{align*}
    [a x^n] [b x^m] &\reducesto [ax^{n+r}bx^{m-r}]\\
    [a x^nb x^m] &\reducesto [ax^{n+r}][bx^{m-r}].
  \end{align*}
\end{lemma}

\begin{figure}
\centering{
\resizebox{60mm}{!}{\svgfig{abcrossings}}
\caption{}
\label{fig:abcrossings}
}
\end{figure}

\begin{figure}
\centering{
\resizebox{60mm}{!}{\svgfig{case1}}
\caption{}
\label{fig:abcrossings}
}
\end{figure}

\begin{figure}
\centering{
\resizebox{60mm}{!}{\svgfig{case2}}
\caption{}
\label{fig:abcrossings}
}
\end{figure}


\begin{figure}
\centering{
\resizebox{60mm}{!}{\svgfig{ablifts}}
\caption{}
\label{fig:abcrossings}
}
\end{figure}

\begin{figure}
\centering{
\resizebox{100mm}{!}{\svgfig{abliftsbands}}
\caption{}
\label{fig:abcrossings}
}
\end{figure}



\begin{proof}
  Use the band model, centered on~$x$.
  Number the windows on top and bottom (cut out by other lifts of~$x$)
  with an arbitrary basepoint on top, and chosen so that lifts of $a$ connect
  top window~$\ell$ to bottom window~$\ell$. Choose $k$ so that lifts
  of~$b$ connect top window~$\ell$ to bottom window $\ell+k$.
  See Figure~\ref{fig:abliftsfirst}.
  
  \begin{figure}
\centering{
\resizebox{100mm}{!}{\svgfig{abliftsfirst}}
\caption{}
\label{fig:abliftsfirst}
}
\end{figure}

  For the first reduction, note
  lifts of $a x^n$ connect windows with a shift of~$n$, and
  similarly lifts of $b x^m$ connect windows with a shift of $k+m$.
  Choose a lift of $a x^n$ that connects window $0$ to window~$n$,
  and choose a lift of $b x^m$ that connects window $n+r-k-m$ to
  window $n+r$. Smoothing gives the desired reduction.

  Other reduction is similar, instead choosing two different lifts of
  $a x^n b x^m$.
  
  See DPT notes \texttt{2018-03-01 \#1} (but it's not written up very
  cleanly there).
\end{proof}

\begin{proposition}\label{prop:join-split-alg}
  If $a$ and $b$ cross $x$ to the right, then
  \[
    \lim_{n,m\to \infty} v(a x^n b x^m) - v(ax^n) - v(bx^m) = 0.
  \]
\end{proposition}
\begin{proof}
  Let $k$ be the value from Lemma~\ref{lem:join-split-geom} for $a$ and
  $b$. Then, by that Lemma,
  \[
    v(a x^n) + v(b x^{n+k+4}) \ge v(a x^{n+1} b x^{n+k+3})
      \ge v(a x^{n+2}) + v(b x^{n+k+2}).
  \]
  In particular,
  $\lim_{n \to \infty} v(a x^n b x^{n+k+2}) - v(x^{2n+k+2})$ exists
  (and equals what it is supposed to equal).

  Muck around a little to show that the limit exists in general.
\end{proof}

\subsubsection{$\mu_{\nu}$ is well-defined}

\begin{proposition}
  The limit defining $\mu_{v}(B_{a,b,c,d,x})$ exists and is positive.
\end{proposition}

\begin{proof}
  Existence is clear from Corollary~\ref{cor:additivity-limit}.

  For positivity, show that $[b x^n c x^m] \reducesto [a x^m][d x^n]$.
  Apply Proposition~\ref{prop:join-split-alg}.
\end{proof}

\subsubsection{$\mu_{\nu}$ is side-by-side additive}

We will address additivity of $\mu_v$ more later, but for now we prove
one easy case.

\begin{proposition}
  If two boxes $B_1$ and $B_2$ fit together side-by-side, then
  $\mu_v(B_1 \cup B_2) = \mu_v(B_1) + \mu_v(B_2)$.
\end{proposition}

\begin{proof}
  Define the boxes by elements $p_1,p_2,q_1,q_2,q_3 \in \pi_1(S, *)$.
  Then
  \begin{align*}
    B_1 &= (q_1 p_1, q_2 p_1, q_1 p_2, q_2 p_2)\\
    B_2 &= (q_2 p_1, q_3 p_1, q_2 p_2, q_3 p_2)\\
    B_1 \cup B_2 &= (q_1 p_1, q_3 p_1, q_1 p_2, q_3 p_2).
  \end{align*}
  Additivity follows immediately from the formula (see Figure~\ref{fig:finiteadditivity}).
  
  \begin{figure}
\centering{
\resizebox{120mm}{!}{\svgfig{finiteadditivity}}
\caption{Illustration of side-by-side additivity for boxes}
\label{fig:finiteadditivity}
}
\end{figure}
\end{proof}


\begin{proposition}
The right handed rectangles $B_{a,b,c,d,x}$ form a semi-ring.
\end{proposition}

(Proof is showing more: that we only need side-by-side additivity.
What's the term for that?)

\begin{proof}
Intersection of boxes is a box.
If $A,B$ are boxes, $A - B$ is a union of boxes. Idea: subdivide and induction on the number of vertical+horizontal
coordinates that appear. Look at any vertical division whose whisker
reaches
the upper boundary. If there isn't one, there's a rectangle that
reaches clear across the top; remove that rectangle, by induction
build up the rest using side-by-side additivity, and then add that
missing rectangle.

Otherwise, divide the entire collection of rectangles at the location
of that vertical division. (This is possible, since this is the
correct direction for the whisker.) By induction, build up each half
by side-by-side additivity, and then attach the two sides together.
\end{proof}

Let $\mathcal{R}$ be the ring generated by $\mathcal{S}$.
By Propositon~\ref{prop:addsemitoring}, $\mu_{\nu}$ is a finitely additive function.


\subsection{$\mu_{\nu}$ is countable additivity}
\label{sec:countable-additivity}

We prove countable additivity of the function $\mu_{\nu}$ by constructing a final family.


\begin{proposition}[Final family]

For $a,b,e$ crossing $x$ to the right, define a family
\[
  F_{a,b,e} = \bigl(B_{a,bx^ke,c,dx^k e}\bigr)_{k=0}^\infty
\]

See Figure~\ref{fig:finalfamily} for an illustration.
\begin{figure}
\centering{
\resizebox{120mm}{!}{\svgfig{finalfamily}}
\caption{Box representing an element of the final family $F_{a,b,e}$. Here $e=Ab=Cd$.}
\label{fig:finalfamily}
}
\end{figure}

\begin{itemize}
\item $F_{a,b,e}$ is nested
\item $\lim_{k \to \infty} \mu_v(F_{a,b,e}(k)) = 0$.
\item $F_{a,b,e}$ is final.
\end{itemize}
\label{prop:finalfamily}
\end{proposition}
\begin{proof}
We prove each one of the items:
\begin{itemize}
    \item Obvious by construction.
    \item 
    We want to show


\[
\lim_{k \to \infty} \nu(F_{a,b,e}(k))=\lim_{k \to \infty} \lim_{n,m \to \infty} \nu(ax^kex^ncx^m)-\nu(ax^n)-\nu(cx^kex^m)=0
\]
    
    By Proposition~\ref{prop:join-split-alg}, the limit in question is equivalent to the same limit for the expression
    
   \begin{align*}
       \nu(ax^k) - \nu(x^k) + & \nu(ex^n) - \nu(x^n) + \nu(cx^m) - \nu(x^m) \\
       &- (\nu(ax^n) - \nu(x^n)) - (\nu(cx^k)-\nu(x^k)) - (\nu(ex^m)-\nu(x^m))
    \end{align*}
    Since the order of limits doesn't matter, the claim follows.
    
    
    \item Suppose $(R_i)_{i=1}^{\infty}$ is a nested sequence of elements of $\mathcal{R}$ with empty intersection. By definition of ring, for each $i$, we can write
    
    \[
        R_i = \cup_{j=1}^{n_i} B^i_{a_j,b_j,c_j,d_j}
    \]
    where the sets in the union are pairwise disjoint.
    Since we are assuming the intersection of the $R_i$ is empty, and $B^i_{a_j,b_j,c_j,d_j}$ are products of half-open intervals, the boxes $B^i_j=[p^i_j,q^i_j) \times [r^i_j,s^i_j)$ must be focusing in the open endpoints. Therefore, there exists some $N$ so that for all $i \geq N$, $p_i<p_{i+1}$ and $r_i<r_{i+1}$.
    Consider a choice of $a,b,e$ so that the two open endpoints of $F_{a,b,e}$ are fixed at $q_N,s_N$ and one of the closed endpoints is fixed to $r_N$. We let the other closed endpoint $p_N$ vary.  See Figure~\ref{fig:nestedhalfopen}.
    We can write $F_{a,b,e}(k)$ as a box of type $[p_k,q_{k}) \times [r_{k},s_{k})$. At each level $i$, the boxes $B^i_j$ focus on open endpoints. Thus, there exists $i$ so that for all $j$, $p^i_j>p_k$ and $r^i_j >r_k$.
     Let $X = \bigcap_i \overline{R_i}$, so that $X$ is non-empty and
compact.

For each $x \in X$, there is some $i$ so that $x \in \overline{R_i} - R_i$,
i.e., it is on the left or bottom edge of some box and not contained
in anything else in $R_i$. Then there is an open neighborhood $U_x$ of $x$
so that $X \cap U_x$ is contained in a horizontal or vertical line.

By compactness, finitely many such $U_x$ cover $X$, and thus $X$ is
contained in finitely many horizontal or vertical lines.

It thus suffices to consider a finite union of final boxes $(F^i_{a,b,e})_{j=1}^N$ where $(F^i_{a,b,e})_{j+1} \subset (F^i_{a,b,e})_j$.


\begin{figure}
\centering{
\resizebox{100mm}{!}{\svgfig{finiteadditivity}}
\caption{Finite additivity whiskers.}
\label{fig:limitingbox}
}
\end{figure}

\begin{figure}
\centering{
\resizebox{100mm}{!}{\svgfig{countable_additivity_whiskers}}
\caption{Countable}
\label{fig:limitingbox}
}
\end{figure}
    
\begin{figure}
\centering{
\resizebox{50mm}{!}{\svgfig{nestedhalfopen}}
\caption{Nested half-open boxes.}
\label{fig:nestedhalfopen}
}
\end{figure}


    \end{itemize}
\end{proof}
Sketch:
\begin{enumerate}
\item  By Theorem~\ref{thm:caratheodory}, $\mu_{\nu}$ is a measure on the $\sigma$-ring of right-handed boxes $B_{x,a,b,c,d}$
\item Boxes have dense endpoints, so any other box of geodesics can be approximated.
\item From this follows outer and inner regularity, also locally finite is by construction, so $\mu_{\nu}$ is a Radon measure.
\item By construction $\mu_{\nu}$ is $\pi_1(S,*)$-invariant.
\item $\mu_{\nu}$ is therefore a geodesic current.
\end{enumerate}

\subsection{$\mu_{\nu}$ recovers $\nu$ on curves}
\label{sec:recovery}

\begin{proposition}
For $C$ a closed curve, we have
\[
i(\mu_{\nu},C)=\nu(C).
\]
\end{proposition}
\begin{proof}

\end{proof}

\begin{remark}
This implies in particular that $\nu$ is non-negative. We want measure to be zero on empty set.
\end{remark}

\begin{comment}
\subsection{The proof, older version}
\DMGtodo[inline]{This is the old version of the proof}
\begin{proposition}\label{prop:premeasurewelldef}
$\mu_{v}(B_{a,b,c,d,x})$ is well-defined.
\end{proposition}
\begin{proof}
For the elements $x,a,b,c,d \in \pi_1(S)$, we consider the following curves

\textbf{Type 1}

\begin{enumerate}
\item $cx^ndx^m$
\item $cx^n$
\item $dx^m$
\end{enumerate}

\textbf{Type 2}

\begin{enumerate}
\item $ax^nbx^m$
\item $ax^n$
\item $bx^m$
\end{enumerate}


If we define $f(n)\coloneqq v(cx^n)$, $g(m)\coloneqq v(dx^m)$ and $h(n,m)\coloneqq v(cx^ndx^m)$.
Application of smoothing and additive union yields

\begin{enumerate}
\item $f(n) \sim K_1 + nv(x)$
\item $g(m) \sim K_2 + mv(x)$.
\end{enumerate}

where $\sim$ here means asymptotically.
(show intersections are essential by working in the band model).\\
We want to show that

$h(n,m) \sim K_1 + K_2 + (n+m)v(x)$.

To do this, we use the smoothing property again, and take lifts in the bands model to see that intersections are essential and thus smoothing yields
\[
  [cx^n][dx^{n+k}] \searrow [cx^ldx^{n+k-l}]
\]
where $1 \leq l \leq k-1$.

This shows that
\[
  f(n+k) + g(n) \geq h(n+l,n+k-l).
\]
Similarly,
\[
  [cx^ndx^m][cx^k] \searrow [dx^{n+m-k}]
\]
which yields
\[
  h(n,m) + f(k) \geq g(n+m-k)
\]
i.e.,
\[
  h(n,m) \geq K_1 + K_2 + (n+m)\nu(x) - C(k).
\]
For, $n+m \gg k$ (or $n+m > k + C$). \\

Now we show that $v(cx^n)$ is convex as a function of $n$.
\end{proof}
\end{comment}






%%% Local Variables:
%%% mode: latex
%%% TeX-master: "Smoothings"
%%% End:
